\section{投资逻辑}

渝农商行是本报告排序第一的标的，因为它不是单纯的低估值银行，而是2026Q1已经出现收入和息差改善，同时价格与资金行为尚未形成明显趋势扩散。2026Q1营业收入78.30亿元，同比增长8.39\%；归母净利润39.35亿元，同比增长5.07\%。净利息收入同比增长15.08\%，单季净息差达到1.69\%，较2025年末提升9个基点。对于银行而言，这比单纯的净利润增长更重要，因为它说明负债端重定价开始转化为核心收入改善。

\section{资产质量与增长质量}

\begin{exhibitbox}[Exhibit 6：渝农商行2026Q1关键指标]
\begin{tabularx}{\textwidth}{L{3.6cm}R{2.5cm}X}
\toprule
\textbf{指标} & \textbf{2026Q1} & \textbf{解读} \\
\midrule
营业收入 & 78.30亿元/+8.39\% & 收入增速明显高于多数银行 \\
归母净利润 & 39.35亿元/+5.07\% & 利润增速低于收入，主要因拨备更审慎 \\
净息差 & 1.69\% & 较2025年末提升9bp，核心拐点 \\
不良贷款率 & 1.07\% & 较年初下降1bp \\
拨备覆盖率 & 约366\% & 风险抵补能力仍强 \\
成本收入比 & 约22.4\% & 同比下降约3.1个百分点 \\
\bottomrule
\end{tabularx}
\sourcenote{公司2026年一季度报告；华创证券一季报点评。}
\end{exhibitbox}

风险并未消失。零售贷款需求仍弱，其他非息收入受投资收益波动拖累；若后续通过降低拨备维持利润，质量会下降。但当前收入、息差和不良率同向改善，使其比一般“高股息低增长”银行多了一层基本面修复逻辑。

\section{资金结构}

截至7月10日，渝农商行连续8日主力净流入1.46亿元，同期上涨3.11\%；7月10日股价下跌1.10\%时，公开资金口径仍显示净流入约0.69亿元。2026Q1末沪股通名义持股约3.61亿股，占总股本3.18\%；另有红利低波ETF进入前十大。融资余额约3.02亿元，占流通市值约0.54\%，说明行情并非高杠杆推动。

\section{估值与动作}

开源证券预计2026E营业收入约292.9亿元、归母净利润128.8亿元、EPS 1.13元、BVPS 12.78元。AStock基准使用0.64倍PB，对应8.18元；叠加市场锚和华创8.82元外部目标，最终目标8.01元，相对6.30元空间约27.1\%。

\begin{itemize}
  \item \textbf{配置区}：6.02--6.35元，适合分批；不追单日放量。
  \item \textbf{确认区}：站稳6.70--7.00元，同时下一季净息差保持1.60\%以上。
  \item \textbf{失效条件}：不良率升破1.15\%、净息差跌破1.55\%，或收入增速重新低于2\%。
\end{itemize}
